\chapter[MESSAGES, BULLETINS ET ANNONCES]{Messages, bulletins et
annonces}\label{messages-bulletins-et-annonces}

Les messages, bulletins et annonces APRS sont des paquets contenant des
chaînes de texte libre, destinés à véhiculer de l'information lisible
par un humain. Un message vise un destinataire spécifique et attend en
général un accusé. Bulletins et annonces visent plusieurs destinataires
et ne sont pas acquittés.

\section{Messages}\label{messages}

Un message est un texte avec un destinataire spécifié. Le destinataire
est un champ fixe 9 caractères (padding avec espaces si besoin) suivant
le DTI \texttt{:}. Le destinataire est suivi d'un autre \texttt{:}, puis
du texte.

Longueur maximale : 67 caractères, ASCII imprimable sauf
\texttt{\textbar{}}, \texttt{\textasciitilde{}}, \texttt{\{}.

Identifiant de message optionnel appendu au texte : \texttt{\{} + 1 à 5
caractères alphanumériques sans espace. Les messages sans identifiant ne
sont pas à acquitter. Les messages avec identifiant sont acquittés par
le destinataire --- l'émetteur retransmet jusqu'à ACK, annulation ou
timeout.

\begin{center}
\begin{tabular}{|c|c|c|c|c|c|}
\hline
\textbf{Champ} & \lit{:} & Addressee & \lit{:} & Message Text & \lit{\{} + Message No \\
\hline
\textbf{Octets} & 1 & 9 & 1 & 0-67 & 1 + 1-5 (optionnel) \\
\hline
\end{tabular}
\end{center}

Exemples : - \texttt{:WU2Z\ \ \ \ \ :Testing} --- message à WU2Z, pas
d'ACK attendu. - \texttt{:WU2Z\ \ \ \ \ :Testing\{003} --- même message,
ACK attendu. - \texttt{:EMAIL\ \ \ \ :msproul@ap.org\ Test\ email} ---
exemple e-mail. APRS ne supporte pas l'e-mail nativement.

\section{ACK de message}\label{ack-de-message}

Format proche du message, mais le texte est \texttt{ack} suivi du numéro
de message.

Ex : \texttt{:KB2ICI-14:ack003}.

\section{Rejet de message}\label{rejet-de-message}

Si une station ne peut pas accepter un message : \texttt{rej} au lieu de
\texttt{ack}.

Ex : \texttt{:KB2ICI-14:rej003}.

\section{ACK multiples}\label{ack-multiples}

Si un message est reçu plusieurs fois, on ACK chaque occurrence. Sur
chemin long, plusieurs copies d'ACK peuvent améliorer les chances de
succès. Espacer les ACK multiples d'au moins 30 s (les digipeaters
suppriment les duplicatas dans un intervalle 30 s).

\section{Groupes de messages}\label{groupes-de-messages}

Une station peut définir des Message Groups (listes d'indicatifs à
écouter en plus de son propre indicatif). Groupes définis localement à
la station réceptrice, servent de filtre.

Adresses lues par tous : \lit{ALL}, \lit{QST}, \lit{CQ}. La
station n'acquitte que les messages qui la ciblent directement, jamais
les messages de groupe. Cela vaut aussi en Alternate Net.

\section{Bulletins généraux}\label{bulletins-guxe9nuxe9raux}

Addressee = \texttt{BLN} + un chiffre + 5 espaces. Émis quelques fois
par heure pendant quelques heures. Info temps-sensible. Longueur max 67
caractères, ASCII imprimable sauf \texttt{\textbar{}}
\texttt{\textasciitilde{}}.

Ex : \texttt{:BLN3\ \ \ \ \ :Snow\ expected\ in\ Tampa\ RSN}

\section{Annonces}\label{annonces}

Similaires aux bulletins, mais avec une lettre majuscule au lieu du
chiffre. Beaucoup moins fréquentes (mais possiblement sur plusieurs
jours). Usuellement pas critiques dans le temps. Typiquement en amont
d'un événement, tandis que bulletins servent pendant l'événement.

Les utilisateurs doivent être alertés à l'arrivée d'un nouveau bulletin
ou annonce.

Ex :
\texttt{:BLNQ\ \ \ \ \ :Mt\ St\ Helen\ digi\ will\ be\ QRT\ this\ weekend}

\section{Bulletins de groupe}\label{bulletins-de-groupe}

Adresse = \texttt{BLN} + un chiffre + nom du groupe (jusqu'à 5
caractères + padding espaces).

Ex : \texttt{:BLN4WX\ \ \ :Stand\ by\ your\ snowplows}

Une station peut lister des groupes d'intérêt. Si liste non vide, seuls
les bulletins de ces groupes (+ bulletins généraux) génèrent des
alertes.

\section{Bulletins du National Weather
Service}\label{bulletins-du-national-weather-service}

Format spécial adressé à \texttt{NWS-xxxxx} pour signaler des zones de
la carte concernées par des alertes météo (US).

Ex : \texttt{:NWS-WARN\ :092010z,THUNDER\_STORM,AR\_ASHLEY,\{S9JbA}

L'« identifiant de message » \texttt{\{S9JbA} n'est qu'une référence ---
les bulletins ne sont pas acquittés.

\section{Radiogrammes NTS}\label{radiogrammes-nts}

APRS peut transporter des radiogrammes NTS via le format message
standard, en préfixant le texte par un identifiant
\texttt{Nx\textbackslash{}} :

\begin{itemize}
\tightlist
\item
  \texttt{N\#\textbackslash{}number\textbackslash{}precedence\textbackslash{}handling\textbackslash{}originator\textbackslash{}check\textbackslash{}place\textbackslash{}time\textbackslash{}date}
\item
  \texttt{NA\textbackslash{}address\_line1\textbackslash{}...\textbackslash{}address\_line4}
\item
  \texttt{NP\textbackslash{}phone\ number}
\item
  \texttt{N1\textbackslash{}} à \texttt{N6\textbackslash{}} : lignes 1 à
  6 du texte NTS
\item
  \texttt{NS\textbackslash{}signature}
\item
  \texttt{NR\textbackslash{}Received\ from\textbackslash{}date\_time\textbackslash{}sent\_to\textbackslash{}date\_time}
\end{itemize}

Format inspiré du formulaire ARRL NTS Radiogram. Chaque ligne max 67
caractères y compris l'identifiant 3 caractères. Max 6 lignes de texte
NTS. Champs requis : \texttt{N\#\textbackslash{}},
\texttt{NA\textbackslash{}}, \texttt{NS\textbackslash{}},
\texttt{NR\textbackslash{}}.

Séparateur backslash pour permettre l'inclusion de forward slashes dans
les messages (le backslash n'existe ni en RTTY ni en CW). Sérialisation
par le Message ID \texttt{\{xxxxx} standard.

Une application APRS n'est pas tenue de comprendre ces messages --- ils
restent lisibles en affichage message.

\section{Format obsolète}\label{format-obsoluxe8te}

Certaines stations émettent bulletins et annonces sans le DTI
\texttt{:}. Format obsolète. À interpréter comme Status Report.

\section{Recommandations
d'implémentation}\label{recommandations-dimpluxe9mentation}

Les bulletins/annonces sont un « panneau d'affichage » commun visible
par tous les participants. Info à afficher, à mettre à jour, à retirer.
Espace limité. À implémenter comme « Bulletin Screen » commun.

Points à traiter : - \textbf{Tri} par indicatif émetteur et
Bulletin/Announcement ID. - \textbf{Remplacement} : nouvelles lignes
émises avec même ID écrasent les anciennes de la même station. -
\textbf{Effacement} : l'utilisateur peut retirer une ligne lue. Les
bulletins toujours valides reviendront par retransmission. -
\textbf{Alertes} : sonnerie répétée jusqu'à acknowledgement utilisateur.
- \textbf{Journalisation} : anciens bulletins archivés séquentiellement.

